\chapter{符号测度}
\label{chap:signed-measures}

递增函数的增量可以看成区间所携带的质量；BV 函数则允许正负增量。
另一方面，积分 \(\int_Eg\,d\mu\) 随集合 \(E\) 变化，也具有可数可加性。
本章将这两种来源放在同一种对象中，并回答何种测度能够写成“密度乘以
原测度”。这一问题的答案还将用于描述 \(L^p\) 上的连续线性泛函。

\section{有限符号测度}

固定可测空间 \((X,\mathcal A)\)。本章首先处理有限对象，以免在减法中
出现 \(\infty-\infty\)。

\begin{definition}
有限符号测度是可数可加映射 \(\nu:\mathcal A\to\R\)，并满足
\[
 \sup_{\{E_j\}}\sum_j|\nu(E_j)|<\infty,
\]
其中上确界遍历 \(X\) 的全部有限可测分割。这个有界性称为有限变差条件。
\end{definition}

若 \(\nu=\alpha-\beta\)，其中 \(\alpha,\beta\) 是有限正测度，则上述条件
自动成立。下面证明逆命题，并且可使正负质量集中在互不相交的集合上。
有限变差保证所有不交可测集合列的值绝对可和：任意有限部分的绝对值和
均不超过定义中的上界。若 \(A_n\uparrow A\)，把 \(A\) 分成
\(A_1,A_2\setminus A_1,\ldots\)，可数可加性给出
\(\nu(A_n)\to\nu(A)\)。对 \(A_n\downarrow A\)，应用这一结论于
\(A_1\setminus A_n\)，同样得到 \(\nu(A_n)\to\nu(A)\)。

两个正测度称为互相奇异，若存在可测分割 \(X=P\cup N\)，使一个在
\(N\) 上为零，另一个在 \(P\) 上为零；也就是说，它们分别集中在两个
不交的可测集合上。

\begin{theorem}[Hahn 分解与 Jordan 分解]\label{thm:hahn-jordan-measure}
存在可测分割 \(X=P\cup N\)，使 \(P\) 的每个可测子集上 \(\nu\ge0\)，
\(N\) 的每个可测子集上 \(\nu\le0\)。令
\(\nu^+(E)=\nu(E\cap P)\)、\(\nu^-(E)=-\nu(E\cap N)\)，则
\(\nu=\nu^+-\nu^-\)，且这两个互相奇异的有限正测度唯一。
\end{theorem}
\begin{proof}
令 \(\alpha=\inf_{E\in\mathcal A}\nu(E)>-\infty\)，选
\(\nu(E_n)<\alpha+2^{-n}\)，其中 \(n\ge1\)。由
\[
 \nu(A\cup B)=\nu(A)+\nu(B)-\nu(A\cap B)
 \le\nu(A)+\nu(B)-\alpha
\]
归纳并取递增极限，\(F_n=\bigcup_{k\ge n}E_k\) 满足
\(\alpha\le\nu(F_n)\le\alpha+\sum_{k\ge n}2^{-k}\)。
令 \(N=\bigcap_nF_n\)，连续性给 \(\nu(N)=\alpha\)。
若 \(B\in\mathcal A\)、\(B\subset N\)，则 \(\nu(N\setminus B)\ge\alpha=\nu(N)\)，故
\(\nu(B)\le0\)。若可测集 \(B\subset P=X\setminus N\)，则
\(\nu(N\cup B)\ge\alpha\)，故 \(\nu(B)\ge0\)。

所定义的正负部分是正测度。对任意另一表示 \(\nu=\alpha_1-\beta_1\)，
有 \(\nu^+(E)=\nu(E\cap P)\le\alpha_1(E)\)，同样
\(\nu^-(E)\le\beta_1(E)\)。若另一表示也互相奇异，取其不交集中集
\(P_1,N_1\)，则 \(\alpha_1(E)=\nu(E\cap P_1)\le\nu^+(E)\)，
\(\beta_1(E)=-\nu(E\cap N_1)\le\nu^-(E)\)。两向不等式给出唯一性。
\end{proof}

定义全变差测度 \(|\nu|=\nu^++\nu^-\)。对每个 \(E\)，它等于
\(E\) 的有限可测分割上 \(\sum|\nu(E_j)|\) 的上确界：上界由三角不等式，
下界由分成 \(E\cap P,E\cap N\) 取得。若 \(h\in L^1(|\nu|)\)，定义
\(\int h\,d\nu=\int h\,d\nu^+-\int h\,d\nu^-\)，并有
\(|\int h\,d\nu|\le\int|h|\,d|\nu|\)。

\section{绝对连续与奇异}

对正测度 \(\mu\)，记 \(\nu\ll\mu\)，若 \(\mu(E)=0\) 总蕴含 \(\nu(E)=0\)。
对符号测度，这等价于 \(|\nu|\ll\mu\)：零测集的每个可测子集也零测，
由全变差的分割描述即可看出。
记 \(\nu\perp\mu\)，若存在 \(N\in\mathcal A\) 使 \(\mu(N)=0\) 且
\(|\nu|(X\setminus N)=0\)。例如实线上的单点质量 \(\delta_c\) 与长度测度奇异，
而 \(g\,m\) 在 \(g\in L^1(m)\) 时绝对连续于 \(m\)。

\section{Radon--Nikodym 定理}

\begin{theorem}\label{thm:radon-nikodym}
设 \(\mu\) 为 \(\sigma\)-有限正测度，\(\nu\) 为有限符号测度且 \(\nu\ll\mu\)。
则存在 \(g\in L^1(\mu)\)，使
\[
 \nu(E)=\int_Eg\,d\mu\quad(E\in\mathcal A).
\]
\(g\) 在几乎处处意义下唯一，记 \(g=d\nu/d\mu\)。
若同一可测空间上的 \(\nu\) 为 \(\sigma\)-有限正测度且仍满足
\(\nu\ll\mu\)，同样存在非负密度，积分允许为无穷。
\end{theorem}
\begin{proof}
先设 \(\mu,\nu\) 都为有限正测度。考察所有非负可测函数 \(h\)，满足
\(\int_Eh\,d\mu\le\nu(E)\) 对每个 \(E\) 成立。这个类非空，并对取最大值封闭：
把 \(E\) 分成两个函数谁较大的部分即可验证。
令 \(s\) 为此类中 \(\int h\,d\mu\) 的上确界，选积分趋于 \(s\) 的函数列，
用逐项最大值使它递增。单调收敛给出极限 \(g\)，仍属此类且积分等于 \(s\)。
残差 \(\lambda=\nu-g\mu\) 是非负测度，且 \(\lambda\ll\mu\)。

若 \(\lambda(X)>0\)，选 \(\varepsilon>0\) 使
\(\lambda(X)>\varepsilon\mu(X)\)。对符号测度 \(\lambda-\varepsilon\mu\)
作 Hahn 分解，其正集 \(P\) 满足
\((\lambda-\varepsilon\mu)(P)>0\)，从而 \(\mu(P)>0\)，且
\(\varepsilon\mu(E\cap P)\le\lambda(E\cap P)\le\lambda(E)\)。
于是 \(g+\varepsilon\mathbf1_P\) 仍属所考察的类，积分却严格大于 \(s\)，矛盾。
故 \(\lambda=0\)。

若 \(\mu\) 仅 \(\sigma\)-有限，将 \(X\) 分成有限 \(\mu\)-测度的不交块，
在每块上应用有限情形，再拼接密度。若 \(\nu\) 也是 \(\sigma\)-有限正测度，
先取对两种测度都有限的共同可数分块，论证相同。符号情形分别处理
\(\nu^+,\nu^-\)，其密度之差可积，因为两部分总质量有限。
最后证明唯一性。对正测度情形，取对 \(\mu,\nu\) 都有限的可数分块
\(X=\bigcup_k B_k\)。在每块上，\(g,h\) 的积分都等于有限数
\(\nu(B_k)\)，因而二者有限几乎处处。排除这些零测例外后，令
\(E=B_k\cap\{g\ge h+1/n\}\)，其中 \(n\ge1\)。两积分有限，可以相减，得到
\[
 0=\int_Eg\,d\mu-\int_Eh\,d\mu
   =\int_E(g-h)\,d\mu\ge\frac{\mu(E)}n.
\]
因此 \(\mu(E)=0\)。对所有 \(k,n\) 取并，再交换 \(g,h\)，得到
\(g=h\) 几乎处处。有限符号测度的两个候选密度都在 \(L^1(\mu)\) 中，
同一论证可直接用于它们的差，不涉及无穷积分相减。
\end{proof}

证明从“已被一个密度解释的质量”出发取极大值，再说明任何非零残差都能
继续贡献一块正密度。绝对连续假设正用在确认那块集合具有正的
\(\mu\)-测度；没有它，残差可能完全躲在 \(\mu\)-零测集上。

有限复测度可定义为两个有限实符号测度的组合 \(\nu_1+i\nu_2\)。
分别应用定理得到复密度 \(g_1+ig_2\)。其积分按实部、虚部定义；
后文复 \(L^p\) 的对偶表示只需要这个版本。

\section{Lebesgue 分解与链式法则}

先说明怎样把密度公式用于一般被积函数。若正测度满足
\(\nu(E)=\int_Eg\,d\mu\)，其中 \(g\ge0\)，则对任意非负可测函数 \(h\)，
\[
 \int h\,d\nu=\int hg\,d\mu.
\]
对指示函数，这是密度的定义；对非负简单函数，由有限可加性得到；
再取递增简单函数逼近 \(h\)，两侧分别用单调收敛定理即可。
在 \(g=0\) 处将乘积规定为零，因而这一论证也允许 \(h\) 取无穷值。
对可积实函数或复函数，分解正负部分或实虚部得到相同公式。

\begin{theorem}\label{thm:lebesgue-decomposition}
设 \(\mu,\nu\) 为 \(\sigma\)-有限正测度，则有唯一分解
\(\nu=\nu_a+\nu_s\)，其中 \(\nu_a\ll\mu,\ \nu_s\perp\mu\)。
有限符号测度也有同样的唯一分解。
\end{theorem}
\begin{proof}
先设两者有限，令 \(\rho=\mu+\nu\)。Radon--Nikodym 定理给
\(d\mu=f\,d\rho,\ d\nu=g\,d\rho\)，其中 \(f,g\ge0,\ f+g=1\) 几乎处处。
在一个 \(\rho\)-零集上修改后，可取 \(f,g\) 处处有限且非负。
令 \(A=\{f>0\}\)，把 \(g/f\) 在 \(A^c\) 上补零，并令
\(d\nu_a=\mathbf1_A(g/f)\,d\mu\)、\(\nu_s=\nu|_{A^c}\)。由刚证的换密度公式，
\[
 \nu_a(E)=\int_{E\cap A}\frac{g}{f}\,d\mu
          =\int_{E\cap A}g\,d\rho=\nu(E\cap A).
\]
故两者之和为 \(\nu\)，而 \(\mu(A^c)=\int_{A^c}f\,d\rho=0\)。
一般情形用共同有限分块拼接，奇异集中集取各块集中集的可数并。
若有两种分解，除去两者
奇异支撑的并，两个绝对连续部分相等；在该零测并集上它们都为零，因此处处
作为测度相等，奇异部分也相等。符号情形先分解正负部分再相减，唯一性仍同理。
\end{proof}

若 \(\lambda\ll\nu\ll\mu\) 是 \(\sigma\)-有限正测度，则
\[
 \frac{d\lambda}{d\mu}
 =\frac{d\lambda}{d\nu}\frac{d\nu}{d\mu}\quad\mu\text{-几乎处处}.
\]
选择非负密度，并在 \(d\nu/d\mu=0\) 的集合上把右侧规定为零。
对每个可测集 \(E\)，在换密度公式中代入
\(h=\mathbf1_E(d\lambda/d\nu)\)，再由RN密度的唯一性得到结论。
这里的等式是密度的几乎处处等式，不是逐点函数的商法则。

\section{BV 函数的测度增量}

\subsection{右连续代表}

本节固定有限实数 \(a<b\)，\(F:[a,b]\to\R\) 为有限实值函数，
所有测度均定义在 \(((a,b],\mathcal B((a,b]))\) 上，\(m\) 表示长度测度在此的限制。
设 \(F\) 递增且右连续于 \([a,b)\)。在 \((a,b]\) 上存在唯一有限正 Borel
测度 \(\mu_F\)，满足
\[
 \mu_F((s,t])=F(t)-F(s)\quad(a\le s<t\le b).
\]
若 \(F(b)=F(a)\)，取零测度即可。否则，一种直接构造是把区间
\((0,F(b)-F(a)]\) 的长度测度经
\[
 Q(y)=\inf\{t\in[a,b]:F(t)-F(a)\ge y\}
\]
推向 \((a,b]\)。上式取下确界的集合非空，因为包含 \(b\)。记其下确界为 \(q\)。
若 \(q<b\)，从该集合选点趋于 \(q\)，右连续性给出
\(F(q)-F(a)\ge y\)；若 \(q=b\)，则 \(b\) 本来就在集合中。
故下确界总能取得。由单调性，\(Q(y)\le t\) 等价于
\(y\le F(t)-F(a)\)；又因 \(y>0\)，取得下确界还说明 \(Q(y)>a\)。
因此 \(Q\) 可测，所得推前测度具有所需区间值。唯一性由有限测度在生成
Borel \(\sigma\)-代数的半开区间族上的唯一性得到。

若 \(F\) 右连续且 BV，其累计变差也右连续。为核实这一点，固定 \(x<t\)，
取分割 \(x=s_0<s_1<\cdots<s_m=t\)，其变差和大于 \(V_x^t(F)-\eta\)。
当 \(x<y<s_1\) 时，由变差可加性和在 \([y,t]\) 上使用分割
\(y,s_1,\ldots,s_m\)，得
\[
 V_x^y(F)<\eta+|F(s_1)-F(x)|-|F(s_1)-F(y)|
 \le\eta+|F(y)-F(x)|.
\]
先令 \(y\downarrow x\)，再令 \(\eta\downarrow0\)，得到 \(V_x^y(F)\to0\)，
亦即累计变差右连续。故规范 Jordan 部分右连续，取其 Stieltjes 测度之差定义
\(dF\)。端点 \(a\) 不携带质量；改变 \(F(a)\) 可能改变紧邻端点的增量，
所以右连续性在 \(a\) 处也不可省略。

\subsection{导数能记录的部分}

\begin{proposition}\label{prop:bv-measure-ac}
对上述右连续 BV 函数，\(F\in AC[a,b]\) 当且仅当 \(dF\ll m\)。
此时 \(dF=F'\,m\)。一般情况下，\(dF\) 的 Lebesgue 分解为
\[
 dF=F'\,m+\nu_s,\qquad \nu_s\perp m.
\]
\end{proposition}
\begin{proof}
绝对连续时用基本定理在所有半开区间上计算增量，再用测度唯一性。
反之，Radon--Nikodym 定理给 \(dF=g\,m\)，故
\(F(t)-F(a)=\int_a^tg\)，于是绝对连续。

为识别一般分解的密度，先设 \(F\) 递增。单调可导定理在每个子区间上给
\(\int_s^tF'\le F(t)-F(s)\)，所以
\(H(t)=F(t)-\int_a^tF'\) 仍递增，且 \(dH=dF-F'm\) 为正测度。
特别地，若 \(dF\) 奇异，\(0\le F'm\le dF\) 同时绝对连续又奇异，
只能为零，故 \(F'=0\) 几乎处处。
现在对 \(dF=\nu_a+\nu_s\) 的两个正测度分别取累计函数
\(G(t)=\nu_a((a,t])\)、\(S(t)=\nu_s((a,t])\)，并令 \(G(a)=S(a)=0\)。
有限测度的递减连续性说明两函数在 \([a,b)\) 右连续，其半开区间增量
分别等于对应测度的区间值。故刚证的奇异递增情形可以用于 \(S\)。前者由密度积分
构成，导数等于密度；后者奇异，导数几乎处处为零。两者之和与 \(F\)
只差常数，故绝对连续部分的密度为 \(F'\)。一般 BV 情形对两个递增部分相减。
\end{proof}

正测度 \(\mu\) 的一个\emph{原子}是满足 \(\mu(E)>0\) 的可测集 \(E\)，其每个可测子集的测度
不是零就是 \(\mu(E)\)。不存在这种集合时称测度无原子。
阶跃对应单点上的原子质量，Cantor 函数对应无原子的奇异质量（见习题），绝对连续函数对应
密度质量。三者的区分说明：导数为零可能表示函数没有变化，也可能表示变化
不以长度密度的形式出现。

\section*{习题}
\addcontentsline{toc}{section}{习题}
\subsection*{A组}
\begin{enumerate}[label=\textbf{\thechapter.\arabic*.}]
\exerciseitem{signed-density-variation}设 \(g\) 为实值 \(L^1(\mu)\) 函数，\(d\nu=g\,d\mu\)，证明 \(d|\nu|=|g|\,d\mu\)。
\exerciseitem{signed-hahn-nonunique}说明 Hahn 分解的集合未必唯一，而 Jordan 测度唯一。
\exerciseitem{signed-ac-singular-zero}证明同时绝对连续于 \(\mu\) 又与 \(\mu\) 奇异的有限符号测度为零。
\exerciseitem{signed-dirac-decomposition}求 \(\nu=m|_{[0,1]}+\delta_{1/2}-2\delta_2\) 相对于长度测度的分解及全变差。
\end{enumerate}
\subsection*{B组}
\begin{enumerate}[label=\textbf{\thechapter.\arabic*.},resume]
\exerciseitem{rn-uniform-ac}若 \(\mu\) 为 \(\sigma\)-有限且 \(\nu\) 为有限符号测度，证明 \(\nu\ll\mu\) 等价于：任意 \(\varepsilon>0\) 都有 \(\delta>0\)，使 \(\mu(E)<\delta\) 时 \(|\nu|(E)<\varepsilon\)。
\exerciseitem{rn-failure-nonsigma}在同一可测空间 \(([0,1],\mathcal B([0,1]))\) 上，令 \(\mu\) 为计数测度，\(\nu\) 为长度测度。证明 \(\nu\ll\mu\)，但不存在 \(\nu=g\mu\) 的密度。
\exerciseitem{stieltjes-jump-mass}证明 \(\mu_F(\{c\})=F(c)-F(c-)\)，其中 \(a<c\le b\)。
\exerciseitem{cantor-measure-singular}证明 Cantor 函数产生的概率测度奇异于长度测度且没有原子。
\end{enumerate}
\subsection*{C组}
\begin{enumerate}[label=\textbf{\thechapter.\arabic*.},resume]
\projectexerciseitem{signed-variation-partition}若 \(F\) 右连续且 BV，证明 \(|dF|((s,t])=V_s^t(F)\)，并说明为什么右连续代表的要求不能随意删去。
\end{enumerate}
